\element{differentielle}{
  \begin{questionmult}{diffExacte}\bareme{1}
    $\cos(y)dx - x\sin(y)dy$ est
    \begin{multicols}{2}
      \begin{reponses}
        \bonne{une forme différentielle}
        \bonne{une différentielle}
      \end{reponses}
    \end{multicols}
  \end{questionmult}
}
\element{differentielle}{
  \begin{questionmult}{diffExacte}\bareme{1}
    $xdx+x^2dy$ est
    \begin{multicols}{2}
      \begin{reponses}
        \bonne{une forme différentielle}
        \mauvaise{une différentielle}
      \end{reponses}
    \end{multicols}
  \end{questionmult}
}
\setgroupmode{differentielle}{cyclic}

\element{modes}{
  \begin{questionmult}{modesTransport1}\bareme{haut=2,p=0}
    Cocher les affirmations exactes
    \begin{reponses}
      \mauvaise{L'eau est un meilleur conducteur thermique que le béton.}
      \bonne{La chaleur est conduite du chaud vers le froid.}
      \bonne{La chaleur peut être transmise dans un matériau macroscopiquement immobile.}
      \bonne{La chaleur peut être transmise dans le vide.}
      \bonne{L'équation de la chaleur sans terme source est irréversible.}
    \end{reponses}
  \end{questionmult}
}
\element{modes}{
  \begin{questionmult}{modesTransport2}\bareme{haut=2,p=0}
    Cocher les affirmations exactes
    \begin{reponses}
      \bonne{Le béton est un meilleur conducteur thermique que l'eau.}
      \bonne{La chaleur est conduite du chaud vers le froid.}
      \bonne{La chaleur peut être transmise dans un matériau macroscopiquement immobile.}
      \bonne{La chaleur peut être transmise dans le vide.}
      \mauvaise{L'équation de la chaleur sans terme source est réversible.}
    \end{reponses}
  \end{questionmult}
}
\setgroupmode{modes}{cyclic}

\element{ppeInfinitesimal}{
  \begin{question}{1erppeinfinitesimal}\bareme{1}
    Le premier principe de la thermodynamique s'écrit, en formulation infinitésimale
    \begin{multicols}{2}
      \begin{reponses}
        \bonne{$dU+dE=\delta W+\delta Q$}
        \mauvaise{$\delta U + \delta E = W + Q$}
        \mauvaise{$U + E = dW + dQ$}
        \mauvaise{$\delta U+\delta E = dW + dQ$}
      \end{reponses}
    \end{multicols}
  \end{question}
}
\element{ppeInfinitesimal}{
  \begin{question}{2eppeinfinitesimal}\bareme{1}
    Le premier second de la thermodynamique s'écrit, en formulation infinitésimale
    \begin{multicols}{2}
      \begin{reponses}
        \bonne{$dS=\delta S_e+\delta S_c$}
        \mauvaise{$\delta S = S_e + S_c$}
        \mauvaise{$S = dS_e + dS_c$}
        \mauvaise{$\delta S = dS_e + dS_c$}
      \end{reponses}
    \end{multicols}
  \end{question}
}
\setgroupmode{ppeInfinitesimal}{cyclic}

\element{diffusion}{
  \begin{question}{equationDiffusion}\bareme{1}
    L'équation de la diffusion thermique s'écrit
    \begin{multicols}{2}
      \begin{reponses}
        \bonne{$\frac{\partial T}{\partial t}-D_{th}\Delta T=0$}
        \mauvaise{$\frac{\partial^2 T}{\partial t^2}+D_{th}\Delta T=0$}
        \mauvaise{$\frac{\partial^2 T}{\partial t^2}-D_{th}\Delta T=0$}
        \mauvaise{$\frac{\partial^2 T}{\partial t^2}-D_{th}\Delta^2 T=0$}
      \end{reponses}
    \end{multicols}
  \end{question}
}

\element{resistance}{
  \begin{question}{resistanceThermique}\bareme{1}
    La résistance thermique d'un barreau cylindrique de section $S$, de longueur $L$ et de conductivité $\lambda$ s'écrit
    \begin{multicols}{4}
      \begin{reponses}
        \bonne{$\frac{L}{\lambda S}$}
        \mauvaise{$\frac{\lambda}{L S}$}
        \mauvaise{$\frac{S}{\lambda L}$}
        \mauvaise{$\frac{L\lambda}{S}$}
      \end{reponses}
    \end{multicols}
  \end{question}
}

\element{diffusivite}{
  \begin{question}{diffusivite}\bareme{1}
    Le coefficient de diffusion a pour expression
    \begin{multicols}{4}
      \begin{reponses}
        \bonne{$\frac{\lambda}{\mu c_V}$}
        \mauvaise{$\frac{\mu}{\lambda c_V}$}
        \mauvaise{$\frac{c_V}{\mu\lambda}$}
        \mauvaise{$\frac{\mu c_V}{\lambda}$}
      \end{reponses}
    \end{multicols}
  \end{question}
}

\element{serieParallele}{
  \begin{questionmult}{serie}\bareme{haut=2,p=0}
    Des résistances thermiques en série
    \begin{reponses}
      \bonne{Sont traversées par le même flux thermique.}
      \mauvaise{Sont soumises à la même différence de température.}
      \bonne{Ont pour résistance équivalente $R_1+R_2$.}
      \mauvaise{Ont pour résistance équivalente $\frac{R_1R_2}{R_1+R_2}$.}
      \mauvaise{Ont pour résistance équivalente $\frac{1}{R_1}+\frac{1}{R_2}$.}
      \bonne{Modélisent le cas d'un mur en béton isolé par de la laine de verre.}
      \mauvaise{Modélisent le cas d'un mur percé d'une fenêtre.}
    \end{reponses}
  \end{questionmult}
}
\element{serieParallele}{
  \begin{questionmult}{parallele}\bareme{haut=2,p=0}
    Des résistances thermiques en parallèle
    \begin{reponses}
      \mauvaise{Sont traversées par le même flux thermique.}
      \bonne{Sont soumises à la même différence de température.}
      \mauvaise{Ont pour résistance équivalente $R_1+R_2$.}
      \bonne{Ont pour résistance équivalente $\frac{R_1R_2}{R_1+R_2}$.}
      \mauvaise{Ont pour résistance équivalente $\frac{1}{R_1}+\frac{1}{R_2}$.}
      \mauvaise{Modélisent le cas d'un mur en béton isolé par de la laine de verre.}
      \bonne{Modélisent le cas d'un mur percé d'une fenêtre.}
    \end{reponses}
  \end{questionmult}
}
\setgroupmode{serieParallele}{cyclic}